\documentclass[class=report, crop=false]{standalone}
\usepackage[T1]{fontenc}
\usepackage[utf8]{inputenc}
\usepackage{amsmath}
\usepackage{amssymb}
\usepackage{bm}
\usepackage[authoryear, round]{natbib}

\begin{document}

\chapter{Literature Review}

\section{Development of Mie Series and its Role in Analytical Solutions}

The theoretical bedrock of electromagnetic scattering is the Mie Series which provides an exact closed form solution to Maxwell's equations for canonical particles by expanding the electromagnetic fields into infinite sums of vector spherical or cylindrical harmonics \citep{Martin1993}. Although \citet{Martin1993} notes that this analytical approach is not a versatile design tool for arbitrary geometries due to its strict requirement that object boundaries must coincide perfectly with the coordinate surfaces of an orthogonal system it remains the absolute reference point for accuracy because it eliminates the discretization and staircasing errors inherent in numerical approximations \citep{Wriedt2012}. Consequently in modern computational electromagnetics the comparison of numerical results against these analytical Mie solutions serves as the primary validation method to verify solver accuracy and ensure the absence of spurious modes before a solver is applied to geometrically complex bodies \citep{Paulsen1994}. To rigorously bridge the gap between analytical theory and full three dimensional numerical implementation \citet{Sjoberg2025} advocates for the use of rotationally symmetric solvers as a critical validation stage. \citet{Sjoberg2025} theoretically justifies this approach by demonstrating that for bodies of revolution the three dimensional vector wave equation can be decomposed into a sequence of decoupled two dimensional scalar problems via azimuthal mode expansion which significantly reduces computational cost while maintaining mathematical exactness. This allows the underlying finite element machinery including the weak form assembly and boundary condition enforcement to be verified against Mie theory with extremely high mesh fidelity before transitioning to computationally expensive arbitrary 3D geometries \citep{Sjoberg2025}. In the specific context of this project the Mie series serves as this analytical benchmark for the scattering of a Perfect Electric Conductor PEC cylinder because the separation of variables in cylindrical coordinates allows for the rigorous derivation of exact scattering coefficients by forcing the tangential electric field to vanish at the conducting surface which provides a direct spectral test for the developed solver.

\section{Comparative Analysis of Numerical Approaches}

The selection between time domain and frequency domain numerical approaches is primarily dictated by the spectral characteristics of the optical problem and the computational efficiency required for steady state analysis. While time domain methods are inherently advantageous for broadband simulations where a single transient pulse can recover the spectral response over a wide frequency range \citep{teixeira2023finite} \citet{Miller1994} argues that this efficiency is strictly inverted for resonant or monochromatic systems because the time domain approach relies on determining the late time response where the signal must be tracked for an extended duration until the transient energy fully dissipates. This makes time stepping algorithmically inefficient for high quality factor structures where the fields ring for thousands of periods compared to frequency domain solvers which directly compute the asymptotic steady state solution without iterating through the entire temporal history \citep{Miller1994}. Furthermore \citet{Anees2019} emphasize that the mathematical formulation of time domain finite element methods requires rigorous discretization in both space and time using complex integration schemes such as backward Euler or symplectic algorithms to ensure numerical stability. This creates a significant computational burden because the accuracy of the temporal evolution is strictly coupled to the spatial mesh density meaning that the fine unstructured meshes required to resolve nanophotonic features necessitate infinitesimally small time steps to avoid instability or excessive dispersion error \citep{Anees2019}. Consequently for the analysis of specific photonic devices operating at fixed laser wavelengths the frequency domain approach is superior as it transforms the hyperbolic wave equation into an elliptic boundary value problem that yields the exact time harmonic solution in a single algebraic step thereby avoiding the cumulative dispersion errors and heavy iteration costs associated with long duration time stepping \citep{Shin2012}.

\vspace{12pt}

The finite-difference frequency-domain (FDFD) method provides a direct algebraic approach to electromagnetic modeling by discretizing the frequency-domain Maxwell's equations, specifically the Helmholtz equation, onto a structured grid \citep{Rumpf2014}. Its implementation typically utilizes the standard Yee cell lattice, where electric and magnetic field components are spatially staggered, allowing differential operators to be approximated by simple finite differences \citep{Rumpf2014}. This formulation assembles the simulation into a large linear system ($Ax=b$) that can be solved directly for the steady-state field distribution \citep{Shin2012}. The primary benefit of FDFD lies in its algorithmic simplicity and its flexibility in modeling general anisotropic materials by averaging tensor properties across grid cells \citep{Rumpf2014}. However, a significant limitation is its reliance on rectangular Cartesian grids, which introduces ``staircasing'' discretization errors when modeling curved boundaries; this can degrade accuracy for complex non-orthogonal geometries unless computationally expensive fine meshes are employed \citep{Shin2012}.

\vspace{12pt}

In contrast to volumetric approaches, the Boundary Element Method (BEM) reduces the dimensionality of the problem by transforming Maxwell's equations into boundary integral equations using Green's functions \citep{Buffa2003}. Instead of discretizing the entire simulation volume, BEM discretizes only the surface of the scatterer into patches, solving for equivalent surface currents rather than volumetric fields \citep{Buffa2003}. A major advantage of this implementation is that the underlying Green's functions naturally satisfy the Silver-Müller radiation condition, ensuring exact modeling of wave propagation in unbounded domains without the need for artificial absorbing boundaries \citep{Buffa2003}. While this ensures high accuracy for homogeneous scatterers, the method produces fully populated (dense) system matrices because every boundary element interacts with every other element, rendering it computationally inefficient for problems involving highly inhomogeneous or nonlinear media compared to sparse methods \citep{Buffa2003}.

\vspace{12pt}

The finite element method (FEM) is a numerical approach that solves the variational form of the vector wave equation by discretizing the computational domain into an unstructured mesh \citep{Jin2015}. Unlike structured grids, this unstructured discretization allows for the precise modeling of curved boundaries without geometric approximation errors and naturally accommodates continuously varying material inhomogeneities by defining properties on an element-by-element basis \citep{Paulsen1994}. However, the method is subject to inherent stability concerns, particularly the accumulation of numerical dispersion errors over large electrical distances which is known as the pollution effect, so it necessitates high mesh densities for high-frequency simulations \citep{Babuska1997}. Furthermore, the formulation requires careful implementation to avoid spurious non-physical solutions and necessitates truncating the unbounded simulation domain, which increases the computational problem size compared to boundary-only methods \citep{Jin2015}. Despite these challenges, the sparsity of the resulting system matrices makes FEM a highly effective tool for analysing complex, heterogeneous structures.

\section{Basis Functions and Spurious Modes}

The selection of basis functions serves as the fundamental component that dictates the physical fidelity of the finite element method in electromagnetic analysis. This choice is not merely numerical but is governed by the strict continuity conditions imposed by Maxwell's equations at the interfaces between different media. In physical reality, the tangential components of the electric field must remain continuous across material boundaries to satisfy Faraday's law, whereas the normal component must undergo a discrete jump proportional to the discontinuity in the electric permittivity. \citet{Bossavit1989} provide a critical justification for rejecting standard node based vectorial finite elements, arguing that they impose continuity on all spatial components of the field across element boundaries. This imposition is a constraint that lacks physical necessity and actively prevents the numerical model from representing the true behavior of the field at dielectric interfaces. When a node-based basis function enforces full continuity in all directions, it artificially smooths the solution, leading to significant errors near boundaries where material properties change abruptly \citep{Bossavit1989}.

\vspace{12pt}

To resolve this conflict between numerical formulation and physical law, the focus shifts to vector basis functions, commonly referred to as edge elements. Unlike their nodal counterparts which associate degrees of freedom with the mesh vertices, these elements assign degrees of freedom to the edges of the computational grid. \citet{Bossavit1989} demonstrate that this geometric association allows the basis functions to define a space where only the tangential continuity is enforced along the element interfaces. This leaves the normal component free to be discontinuous, thereby aligning the mathematical properties of the discretization with the continuity requirements of the electric field \citep{Kamireddy2022}. By ensuring that the numerical subspace reflects the physical space of the electromagnetic field, edge elements allow for the direct modeling of complex internal inhomogeneities without the need for penalty terms or ad hoc smoothing techniques.

\vspace{12pt}

Beyond the issues of interface continuity, the numerical scheme must address the critical issue of spurious modes. These are defined as unphysical solutions that arise within the computational domain, often manifesting as highly divergent fields that pollute the solution spectrum. The origin of these modes lies in the inability of standard nodal elements to properly represent the null space of the curl operator. \citet{Bossavit1989} highlight that edge elements possess the distinct advantage of not generating these spurious modes when applied to resonant cavity problems or scattering formulations. This stability arises because edge elements mimic the exact sequence property of the de Rham complex, ensuring that the gradient of the discrete scalar space is contained exactly within the range of the vector edge element space \citep{Schwarzbach2009}. Consequently, the discrete curl operator mimics the behavior of the continuous operator by mapping gradient fields to exactly zero, effectively separating the unphysical irrotational modes from the physical solenoidal ones.

\vspace{12pt}

While nodal elements combined with potential formulations offer a valid pathway by solving for continuous potentials rather than discontinuous fields \citep{Paulsen1994}, the direct modeling of the electric field using edge elements provides a more robust alignment with the vector nature of Maxwell's equations. This theoretical foundation directly informs the practical development of the finite element solver, which utilizes the open-source platform FEniCSx. The choice of FEniCSx is strategic, as the platform provides a well-established implementation of high order Nédélec elements that are rigorous and mathematically verified. By leveraging these built in capabilities, the solver automatically satisfies the tangential continuity conditions required at material interfaces and possesses the necessary mathematical structure to eliminate spurious modes \citep{Esterhazy2013}. This integration ensures that the simulation tool achieves both physical accuracy and spectral stability in the analysis of arbitrary electromagnetic structures, as strongly justified by the foundational work of \citet{Bossavit1989} and subsequent comparative studies.

\section{Boundary Conditions and Domain Truncation}

The definition of boundary conditions constitutes a pivotal step in the design of a finite element solver, as these constraints serve as the mathematical interface between the finite computational domain and the infinite physical reality of light scattering. The most immediate challenge in scattering simulations is the truncation of the unbounded domain. Since the scattered field propagates outward indefinitely, the numerical model must satisfy the Sommerfeld radiation condition, which requires that the energy radiates away from the source without reflection at infinity \citep{Schwarzbach2009}. To mimic this behavior within a finite mesh, early development stages in this project will prioritize Absorbing Boundary Conditions (ABCs). These apply local differential operators to the boundary surface to approximate the one-way wave equation \citep{Engquist1977}. This choice is driven by the ease of implementation, as ABCs can be integrated directly into the surface integrals of the weak form without requiring the construction of auxiliary geometric domains. However, while computationally inexpensive, these local operators frequently introduce spurious reflections when electromagnetic waves strike the boundary at oblique angles, thereby degrading the accuracy of the solution in the near-field region.

\vspace{12pt}

As the solver development progresses towards higher precision requirements, the focus inevitably shifts from ABCs to the Perfectly Matched Layer (PML). As established by \citet{Berenger1994}, the PML is not a simple boundary condition but rather an artificial absorbing volume surrounding the physical domain. By utilizing complex coordinate stretching or anisotropic material tensors, the PML promotes the exponential decay of electromagnetic waves inside the layer without generating reflections at the interface between the free space and the absorbing medium. \citet{Grand2020-we} provide critical practical justification for this transition, demonstrating that while PMLs are theoretically reflectionless, their implementation requires a sufficiently dense mesh, typically at least five elements across the layer thickness to prevent numerical discretization error. This approach ensures that the solver faithfully emulates an open boundary, allowing for the accurate retrieval of scattering cross-sections and far-field patterns, effectively eliminating the reflection artifacts inherent to simpler ABC implementations.

\vspace{12pt}

Beyond the truncation of exterior boundaries, the imposition of internal boundary conditions is essential for ensuring physical fidelity and exploiting geometric symmetries to reduce computational cost. \citet{Chaber2017} discuss the implementation of Perfect Electric Conductor (PEC) and Perfect Magnetic Conductor (PMC) conditions within the finite element framework. The PEC condition, which forces the tangential electric field to zero, is the standard mathematical representation for highly conductive metals where fields do not penetrate. Conversely, the PMC condition forces the tangential magnetic field to zero. While PMCs do not exist as natural materials, \citet{Chaber2017} highlight their utility in imposing symmetry planes. By analyzing the incident field polarization, a scattering problem on a symmetric object can often be reduced to one-half or one-quarter of the original domain by applying appropriate PEC or PMC conditions on the cut planes, significantly reducing the memory footprint and solution time.

\vspace{12pt}

From the perspective of mathematical stability and formulation, the choice of boundary condition dictates the well-posedness of the system matrix. In the weak variational form used by solvers like FEniCSx, boundary conditions are categorized as either essential (Dirichlet) or natural (Neumann). \citet{Schwarzbach2009} argues that the stability of the finite element solution, particularly at low frequencies, relies on the correct application of these conditions to prevent the singularity of the system matrix. For instance, when solving the curl-curl equation, the lack of proper boundary constraints on the gradient subspace can lead to an ill-conditioned system. By rigorously defining the essential boundary conditions on the electric field (such as PEC on metal surfaces) and allowing the natural boundary conditions to handle the magnetic field continuity via surface integrals, the solver ensures a unique and stable solution. This interplay between physical constraints and mathematical formulation confirms that boundary conditions are not merely auxiliary settings but are intrinsic to the validity of the electromagnetic simulation.

\section{Error Analysis and the Pollution Effect}

The implementation of error analysis within a finite element solver for electromagnetic analysis is not merely a validation step but a theoretical necessity required to certify the physical fidelity of the simulation. In the context of solving time harmonic Maxwell equations, the importance of this analysis stems from the inherent difficulty in distinguishing between a physically converged solution and a numerical artifact caused by discretization errors. \citet{Melenk2024} emphasize that standard error estimation techniques often fail in high frequency regimes due to the pollution effect, where the numerical phase error accumulates as the wave propagates through the mesh. This accumulation leads to a global dispersion error that grows with the wavenumber, meaning that a locally accurate solution can be globally incorrect. Therefore, the primary function of error analysis is to quantify this deviation using specific mathematical norms, ensuring that the chosen discretization strategy possesses the necessary stability to suppress these pollution errors and yield a reliable approximation of the electromagnetic field.

\vspace{12pt}

To perform a rigorous error analysis, one must move beyond simple visual inspection and utilize quantitative metrics established in the mathematical literature. \citet{Zhong2018} provide the theoretical framework for optimal error estimates by utilizing the $H(\text{curl})$ norm alongside the standard $L^2$ norm. The $L^2$ norm measures the average energy error over the domain, while the $H(\text{curl})$ norm accounts for the accuracy of the field's rotation, which is critical for vector electromagnetics. In scenarios where an analytical solution is available, such as the Mie series for a sphere, the error is calculated directly as the difference between the numerical and exact fields. However, for practical engineering problems involving complex geometries where no analytical solution exists, the convergence analysis operates via a process of self-validation \citep{Esterhazy2013}. This involves computing a reference solution on an extremely refined mesh, often referred to as an overkill mesh, and treating it as the ground truth. The error for coarser discretisation is then evaluated relative to this reference solution to verify that the results form a Cauchy sequence, meaning the difference between successive refinements approaches zero \citep{Grand2020-we}.

\vspace{12pt}

The mechanism of convergence analysis in these complex scenarios relies on identifying the asymptotic regime of the solver. \citet{Frelet2024} distinguish between the pre asymptotic regime, where the mesh is too coarse relative to the wavelength causing the error to oscillate unpredictably, and the asymptotic regime, where the error decreases monotonically. By plotting the relative error against the number of degrees of freedom on a logarithmic scale, one can extract the convergence rate. If the solver is functioning correctly, this rate should match the theoretical order of the polynomial basis functions used. Furthermore, for problems lacking a reference solution, one can analyse the residual error, which measures how well the numerical solution satisfies the original partial differential equation locally within each element. Hazard and Lenoir (1996) suggest that comparing the standard curl curl solution against a regularized formulation provides another layer of validation, ensuring that the solution remains stable and free from spurious modes even in the presence of geometric singularities.

\vspace{12pt}

The collective insights from these error and convergence analyses provide a definitive directive for the development of the finite element solver which ensure the adoption of high order elements is indispensable. The theoretical work by \citet{Melenk2024} explicitly demonstrates that mesh refinement alone, known as h refinement, is computationally inefficient for suppressing the pollution effect at high frequencies. They establish a strict stability condition requiring that the ratio of the wavenumber and mesh size to the polynomial degree remains small. This condition implies that increasing the polynomial degree, or p refinement, is the only robust mechanism to relax the mesh density requirements while maintaining accuracy. \citet{Nicaise2020} reinforce this by showing that high order Nédélec elements unlock exponential convergence rates, allowing the solver to reach the stable asymptotic regime with significantly fewer degrees of freedom compared to low order methods. Consequently, the error analysis confirms that utilizing arbitrary order edge elements within the FEniCSx platform is the optimal strategy to ensure both spectral stability and computational efficiency in the simulation of light scattering.

\section{Computational Nanoplasmonics and LSPR Biosensing}

While the foundational principles of finite element modeling apply universally to wave propagation, simulating nanophotonic devices, specifically Localized Surface Plasmon Resonance (LSPR) biosensors, introduces unique physical and computational challenges. LSPR occurs when the electromagnetic field of incident light drives the collective oscillation of conduction electrons in a metallic nanoparticle into resonance \citep{Mayer2011}. This phenomenon concentrates electromagnetic energy tightly at the nanoparticle surface, generating intense near field enhancements. Because the precise resonance condition is exquisitely sensitive to the local refractive index, these nanostructures operate as highly effective transducers for biomedical diagnostics \citep{Das2023}.

\vspace{12pt}

The clinical significance of LSPR technology stems from its capacity for real-time, label-free, and highly sensitive detection. Traditional biomedical assays, such as enzyme-linked immunosorbent assays (ELISA) or fluorescence techniques, necessitate chemical labeling, which can inadvertently alter molecular activity and compromise interaction fidelity. In contrast, LSPR biosensors rely exclusively on monitoring changes in the local refractive index to accurately assess binding properties and trace biomarkers without potential interference. Because these critical interactions occur at the nanoscale, the physical optimization of a nanoparticle's geometric and surface characteristics is essential to maximize signal amplification. Consequently, robust computational simulation becomes a critical prerequisite for optimizing these nanomedical delivery systems and diagnostic devices prior to expensive physical fabrication.

\subsection{Modeling Dispersive Plasmonic Metals}

A primary physical challenge in computational nanoplasmonics is the rigorous treatment of material properties under electromagnetic excitation. This challenge is rooted in the phenomenon of optical dispersion, the physical property where a material's refractive index and relative permittivity vary dynamically as a function of the incident light's frequency or wavelength.

\vspace{12pt}

In standard radio-frequency simulations or basic dielectric wave propagation, background materials such as water or glass are frequently approximated using static, real-valued scalar constants because their optical dispersion is negligible over narrow spectral bands. However, this static assumption completely fails for real plasmonic metals. Gold, the primary transducer in LSPR biosensors, exhibits extreme optical dispersion across the visible and near-infrared spectrum \citep{Mayer2011}. Its relative permittivity must be treated as a complex-valued, wavelength-dependent function: $\epsilon_{\text{Au}}(\lambda) = \epsilon'(\lambda) + i\epsilon''(\lambda)$.

\vspace{12pt}

The physical necessity for this dispersive model lies in the mechanics of plasmon resonance. LSPR is fundamentally triggered by a specific wavelength-matching condition where the real part of the metal's permittivity ($\epsilon'$) becomes strongly negative and precisely offsets the positive permittivity of the surrounding dielectric medium, while the imaginary part ($\epsilon''$) dictates the intensity of optical absorption and ohmic loss \citep{Das2023}. If a static material constant were used, the simulation would be permanently locked to a single state, rendering it mathematically impossible for the solver to observe a resonance peak. Therefore, to capture the true resonant behaviour of the biosensor, it is mandatory to discard static constants and instead incorporate rigorous frequency-dependent empirical data, such as the widely established Johnson and Christy (1972) dataset, which provides exact experimental measurements of gold's complex permittivity.

\vspace{12pt}

While empirical datasets like Johnson and Christy (1972) provide the necessary physical accuracy, they introduce a distinct computational hurdle: they consist of discrete, discontinuously spaced data points, for example measurements taken in intervals of 10 or 20 nanometers. An advanced LSPR simulation, however, aims to predict the bulk sensitivity of a biosensor by tracking resonance peak shifts with sub-nanometer precision. To bridge the gap between discrete experimental data and the need for continuous sub-nanometer resolution, the solver must employ mathematical interpolation techniques, specifically cubic spline interpolation \citep{PhysRevE.100.063305}. Interpolation is a numerical method that constructs a smooth, continuously differentiable piecewise polynomial function that passes exactly through the discrete experimental data points. Rather than forcing the solver to snap to the nearest available experimental frequency, the cubic spline creates a continuous mathematical curve representing the material's behaviour across the entire spectrum.

\vspace{12pt}

This interpolation strategy is the engine that drives the resonance tracking across a broad optical spectrum. By translating the discrete empirical dataset into a continuous mathematical function, interpolation empowers the finite element solver to query the exact complex permittivity at any arbitrary wavelength, for example dynamically calculating $\epsilon_{\text{Au}}$ at exactly 532.4 nm. Consequently, as the solver executes an automated spectral sweep, it can incrementally update the material arrays and free-space wavenumber ($k_0$) at ultra-fine spectral steps. This ensures the solver accurately captures the continuous, smooth evolution of the optical absorption cross-section, allowing for the precise, sub-nanometer identification of the LSPR peak without being artificially constrained by the resolution of the original laboratory dataset.

\subsection{Optical Cross-Sections and the Rayleigh Approximation}

To quantify the LSPR phenomenon computationally, solvers must extract macroscopic optical observables from the raw, microscopic simulated near field. LSPR physically manifests as a sharp, highly localized peak in the attenuation of incident light across a specific wavelength spectrum. This attenuation is formally quantified by the total optical extinction cross-section ($C_{\text{ext}}$), which serves as the primary metric for tracking the LSPR peak \citep{Mayer2011}.

\vspace{12pt}

Total optical extinction is defined as the sum of two fundamental physical processes: the absorption cross-section ($C_{\text{abs}}$) and the scattering cross-section ($C_{\text{sca}}$). Absorption ($C_{\text{abs}}$) represents the electromagnetic energy permanently dissipated as resistive heat (Ohmic loss) within the plasmonic metal. Computationally, it is calculated by volumetrically integrating the simulated electric field intensity ($|E|^2$) against the imaginary part of the metal's complex permittivity over the entire nanoparticle domain. Scattering ($C_{\text{sca}}$) represents the incident energy that is redirected outward into the far field. It is calculated by evaluating the outward power flux via a surface integral of the time-averaged Poynting vector across a closed mesh boundary surrounding the nanoparticle. These cross-sections form the vital mathematical bridge between the finite element solver and real-world experiments. By calculating these metrics, researchers can translate the raw localized electric field vectors into measurable, macroscopic spectral data that exactly mimics the output of a laboratory spectrophotometer.

\vspace{12pt}

However, while computing volumetric absorption is mathematically straightforward within the finite element framework, computing the scattering cross-section presents distinct challenges. In highly parallelized computing environments utilizing the Message Passing Interface (MPI), evaluating continuous surface flux integrals across partitioned internal mesh facets can induce severe numerical instabilities and communication bottlenecks.

\vspace{12pt}

Fortunately, classical scattering theory provides a robust physical simplification. For metallic nanoparticles operating in the deep sub-wavelength regime, specifically particles with radii less than 30 nm, the system enters the Rayleigh scattering limit \citep{Mayer2011}. In this quasistatic regime, optical scattering becomes highly inefficient, and the total extinction is overwhelmingly dominated by internal resistive heating losses. Consequently, the relationship simplifies to $C_{\text{ext}} \approx C_{\text{abs}}$. By leveraging this physical approximation, advanced computational frameworks can safely bypass unstable surface flux integrals. This ensures High-Performance Computing (HPC) stability while preserving the accurate tracking of the LSPR spectral peak required for biosensor validation.

\subsection{Quantifying Biosensor Performance: Bulk Linearity and Decay Rates}

The clinical viability of Localized Surface Plasmon Resonance (LSPR) biosensors is rigorously validated through specific computational performance metrics, primarily bulk sensitivity and the evanescent decay length \citep{Das2023}.

\vspace{12pt}

\noindent\textbf{Bulk Linearity and Sensitivity Calibration}

\vspace{6pt}

The fundamental transduction mechanism of an LSPR sensor dictates that the resonance peak wavelength ($\lambda_{\text{peak}}$) exhibits a shift proportional to variations in the surrounding refractive index ($n$). To computationally validate sensor designs, a systematic bulk linearity analysis must be executed by sweeping the background refractive index across a designated operational range, for example $n = 1.33$ to $1.36$ for aqueous biological media. This parametric sweep facilitates the mathematical derivation of the Bulk Sensitivity ($S$), defined as $S = \Delta\lambda_{\text{peak}} / \Delta n$, and expressed in nanometers per Refractive Index Unit (nm/RIU).

\vspace{12pt}

Establishing a strictly linear response is a critical prerequisite for diagnostic applications \citep{Mayer2011}. In clinical environments, linearity ensures constant sensitivity across diverse analyte concentrations. This allows specific optical shifts to be reliably correlated with unique biological binding events, preventing the degradation of diagnostic accuracy that occurs when complex, non-linear calibration models are applied to varying sample concentrations.

\vspace{12pt}

\noindent\textbf{Evanescent Field Confinement and Decay Rate Analysis}

\vspace{6pt}

Beyond bulk sensitivity, the diagnostic efficacy of an LSPR biosensor is fundamentally dependent on its surface sensitivity, which is governed by the evanescent field. The evanescent field is characterized as a non-propagating electromagnetic wave wherein the energy is highly concentrated at the metal-dielectric interface and decays exponentially into the surrounding background medium \citep{Mayer2011}.

\vspace{12pt}

A major limitation of nanoplasmonic sensors is that their enhanced electromagnetic field only extends a very short distance from the metal surface. This rapid decay creates a strictly limited detection region. If a biological target binds to the sensor but remains outside of this specific boundary, the electromagnetic field cannot interact with it. As a result, the sensor will not produce an optical shift, leaving the binding event completely undetected \citep{Grand2020-we}. Due to this physical distance limit, it is essential to calculate exactly how far the field extends to ensure the sensor's capture layer is designed to hold targets within the active detection range.

\vspace{12pt}

This quantification is achieved computationally through a decay rate analysis. By executing parametric simulations that iteratively vary the physical thickness of the surrounding dielectric matrix, for example a protein sensing shell, and calculating the subsequent peak redshift, the spatial sensitivity profile of the nanoparticle is mapped. The resulting computational data is fitted to an asymptotic exponential model to extract the characteristic decay length ($l_d$). As emphasized by \citet{Das2023} and \citet{Mayer2011}, the extraction of this decay length is critical, as it dictates the maximum functional thickness of the biosensor's capture layer. If the biological receptor layer exceeds $l_d$, the sensor experiences premature saturation, thereby degrading its theoretical Limit of Detection (LOD). Thus, this parametric evaluation confirms that the physical architecture of the sensor is mathematically optimized to maximize the spatial overlap between the electromagnetic enhancement and the target analyte.

\section{Summary and Research Gaps}

The existing literature establishes a comprehensive blueprint for physically accurate electromagnetic solvers by noting that the mathematical framework for two-dimensional frequency-domain finite element methods is well established. Foundational principles such as the use of Nédélec edge elements to ensure tangential field continuity, the application of absorbing boundary conditions for open domain scattering, and rigorous discretization strategies to maintain stability are all extensively documented in current research. However, a critical implementation gap persists that prevents these proven methodologies from being accessible to the broader light scattering research community. This gap manifests across the current landscape of simulation tools by creating significant barriers to entry for many researchers who must choose between prohibitive costs or extreme technical complexity.

The primary challenge in light scattering simulation is not a lack of theoretical knowledge but rather the practical difficulty of accessing or developing reliable computational tools. While established commercial platforms offer powerful tools, their proprietary nature often prevents deep algorithmic verification and limits transparency for independent reproduction. Furthermore, these closed architectures can be difficult to customize for the unique coupled physics scenarios inherent to stimuli-responsive biosensors, such as directly linking mechanical deformation data to optical simulation inputs. Alternatively, raw general-purpose finite element libraries provide high mathematical flexibility but demand deep expertise in computational mathematics and advanced programming. Implementing a complete electromagnetic solver from these first principles requires proficiency in variational formulations, $H(\text{curl})$ function spaces, and complex mesh generation. For biosensor researchers whose primary expertise lies in chemistry, materials science, or experimental optics, these technical requirements represent a significant barrier that often leads to the use of oversimplified analytical models that cannot account for complex deforming geometries. Additionally, a significant portion of electromagnetic research relies on in-house or custom-built codes that are rarely made available to the public and frequently lack the comprehensive documentation, standardized validation suites, and long-term maintenance required for scientific reproducibility. Even when such codes are released, they are often undocumented and require expert-level understanding to adapt, or they may depend on specific proprietary libraries and infrastructure that are not accessible to all researchers.

A critical gap remains in the unification of these advanced methodologies into a single accessible framework. While the literature defines individual components for high-order stability, optimized parallel solvers, and rigorous domain truncation, there is currently no single open-source implementation that efficiently combines high-order edge elements, automated variational assembly, and robust validation protocols specifically tailored for multi-scale scattering problems. This research aims to bridge that gap by integrating these disparate theoretical advancements into a cohesive and verifiable solver that lowers the technical barrier for the light scattering community while maintaining the highest standards of physical and mathematical accuracy.

\end{document}